\documentclass[11pt]{article}
\usepackage{setspace}
\usepackage{geometry}
\geometry{margin=1in}
\setstretch{1.15}

\begin{document}

% --- Applicant Information ---
\noindent
\textbf{Zhongli Wang} \\
Phone: +49 160 9541 6015 \\
Email: nkwangzhongli@gmail.com \\
Website: https://sites.google.com/view/zhongli-wang \\
Bielefeld, Germany \\

\vspace{0.5cm}

% --- Recipient Information ---
\noindent
\textbf{Recruitment Committee} \\
Living Conditions Department \\
Luxembourg Institute of Socio-Economic Research (LISER) \\
Belval, Luxembourg \\

\vspace{0.5cm}

% --- Subject Line ---
\noindent
\textbf{Application for Research Associate (Postdoc) in Applied Econometrics and Distributional Analysis --- Ref. 25-43} \\

\vspace{0.3cm}

\noindent Dear Prof. Andreoli and Members of the Recruitment Committee,\\[0.2cm]

I am writing to apply for the Research Associate (Postdoc) position in applied econometrics and distributional analysis (Ref.\ 25-43) in the Living Conditions Department at LISER. I am currently completing my Ph.D.\ in Economics at Bielefeld University and Universit\`a Cattolica del Sacro Cuore (Marie Sk{\l}odowska--Curie EPOC programme). My research training combines applied microeconometrics and large-scale data work, with a substantive interest in inequality of opportunity and the long-run consequences of policy interventions. The EDO project is therefore an excellent match for both my methodological profile and my motivation to produce policy-relevant causal evidence on education reforms using rich administrative and survey microdata.

My doctoral research centers on identifying behavioral and welfare-relevant impacts of institutional and technological changes under heterogeneity. Methodologically, I have extensive experience building reproducible pipelines for microdata analysis, conducting high-dimensional robustness exercises, and implementing simulation-based estimation when analytical solutions are not available. These skills translate naturally to the EDO agenda, where credible causal identification and careful measurement are central: constructing cohort-by-region exposure measures, validating reform timing and treatment intensity, and estimating heterogeneous effects with modern quasi-experimental methods. In particular, I am comfortable with panel and repeated-cross-section settings, event-study designs, and the practical challenges of working with large linked datasets (documentation, replication, and sensitivity to alternative specifications). I am proficient in Python and R, and I am prepared to work in Stata as required by the project, given my strong experience transferring workflows across statistical environments.

I am particularly excited by EDO's combination of high-quality data infrastructure and clearly motivated identification strategies across the three case studies. For Luxembourg, the opportunity to analyze newly created administrative data on curricula expansion and students' competences and tracking choices is especially compelling. My contribution would focus on (i) transparent construction and validation of analysis files in close collaboration with OEJQS and LISER colleagues, (ii) rigorous model estimation of reform impacts on short-run competences and medium-run educational trajectories, and (iii) distributional and heterogeneity analyses that speak directly to inequality of opportunity (e.g., differential effects by socio-economic background, migrant status, baseline achievement, or school context). I am also keen to contribute to cross-country comparability where feasible, for instance by aligning outcome definitions and harmonizing core robustness checks across Luxembourg, Italy, and Sweden.

Alongside my dissertation, I have coauthored projects that strengthen my fit for EDO through demonstrated strengths in data-intensive empirical work and research transparency. \textit{Polytope Fraud Theory} applies supervised machine learning to detect structured anomalies in firm-level financial statements, reflecting my ability to manage large datasets, engineer features, validate models, and communicate results with interpretability. A second project develops agentic workflows that integrate large-language-model agents into empirical research pipelines, emphasizing reproducibility, documentation, and scalable processing of complex data sources. While EDO is fundamentally a causal applied micro project, these experiences reinforce my capacity to build robust data workflows, automate quality checks, and maintain transparent research practices under sensitive data environments.

Beyond research, I am motivated by LISER's strong applied microeconometrics community and its commitment to policy engagement. I would be glad to contribute to department life (workshops, lecture series, and outreach) and to support the annual Winter School on Inequality and Social Welfare Theory co-organized by CESifo, LISER, and the University of Verona. I value collaborative research settings and would enjoy working closely with the EDO team and international partners to develop publishable outputs in leading peer-reviewed journals, while also producing policy briefs and communication materials where appropriate. I am fluent in English; I can work in multilingual environments and would be happy to strengthen my French for professional purposes.

Thank you so much for considering my application! I would be delighted to contribute to the EDO project and LISER's mission of producing rigorous, policy-relevant evidence on education, inequality, and long-run opportunities. I would welcome the opportunity to discuss my candidacy further.

\vspace{6mm}

\noindent Sincerely,\\[15mm]
\noindent Zhongli Wang

\end{document}